\section{Auswahl und Evaluierung synthetischer Datenerweiterung}
\label{sec:synthetic-data}

\subsection{Inhaltliche Stichpunkte}
\begin{itemize}
  \item \textbf{Vergleichsaufbau aus Notebook~01b}
  \begin{itemize}
    \item Je ein Random-Forest-Modell wird mit Originaldaten, DAE-Daten,
          TabGAN-Daten und dem hybriden Datensatz der Knoten trainiert.
    \item Alle vier Varianten werden laut Notebook auf demselben
          Original-Testdatensatz bewertet.
    \item Macro-F1: Original 0{,}8748; DAE 0{,}7860; TabGAN 0{,}8730;
          Combined Nodes Hybrid 0{,}8890.
    \item ROC-AUC: Original 0{,}9692; DAE 0{,}9327; TabGAN 0{,}9689;
          Combined Nodes Hybrid 0{,}9845.
    \item Test Loss: Original 0{,}2954; DAE 0{,}5892; TabGAN 0{,}2968;
          Combined Nodes Hybrid 0{,}2682.
    \item Die Notebook-Grafik zeigt für den hybriden Datensatz den besten
          Macro-F1- und ROC-AUC-Wert dieser vier Varianten.
  \end{itemize}
  \item \textbf{Klassenabhängige Generatoranteile}
  \begin{itemize}
    \item Vorgesehene Aufteilung der synthetischen Anteile:
          Backdoor 65\,\% TabGAN / 35\,\% DAE;
          none 55\,\% TabGAN / 45\,\% DAE;
          syn-flood 60\,\% TabGAN / 40\,\% DAE.
    \item Die global gezählten Anteile im fertigen Datensatz weichen wegen
          der zusätzlichen Originaldaten ab:
          Backdoor 55{,}78\,\% TabGAN, 30{,}13\,\% DAE,
          14{,}09\,\% Original;
          none 53{,}20\,\% TabGAN, 43{,}51\,\% DAE,
          3{,}29\,\% Original;
          syn-flood 56{,}18\,\% TabGAN, 37{,}48\,\% DAE,
          6{,}35\,\% Original.
    \item Die Pipeline erzeugt 100 Blöcke je Knoten und jeweils
          100\,000 Zeilen pro Knoten (acht Knoten).
  \end{itemize}
  \item \textbf{Begründung der Hybridwahl}
  \begin{itemize}
    \item Notebook-Begründung: TabGAN erhält insgesamt höhere Anteile,
          weil es in den dortigen Modellvergleichen gegenüber DAE höhere
          Macro-F1-/ROC-AUC-Werte und niedrigere Fehlerwerte erzielt.
    \item Für none wird der höhere DAE-Anteil mit der Abbildung von
          Sensorstatistiken sowie Idle-/Charging-Variabilität begründet;
          für Backdoor mit besser erhaltenen Korrelations- und
          Diskriminierungsmustern von TabGAN.
    \item Diese klassenbezogenen Erklärungen anhand zusätzlicher
          Klassenmetriken oder Verteilungsvergleiche belegen; die globale
          Vergleichsgrafik allein zeigt keine klassenweise Generatorgüte.
  \end{itemize}
\end{itemize}

\subsection{Noch zu erklären}
\begin{itemize}
  \item Die Notebook-Spalte \texttt{Reconstruction\_Error} wird mit
        \texttt{mean\_squared\_error} aus kodierten Klassenlabels und
        Vorhersagen berechnet. Das ist kein Rekonstruktionsfehler des
        Generators und sollte entsprechend umbenannt oder nicht als
        Fidelity-Metrik interpretiert werden.
  \item \texttt{validation\_loss} stammt beim Random Forest aus
        Out-of-Bag-Wahrscheinlichkeiten; beschreiben, warum dieser Wert
        nicht mit einem separaten Validierungsdatensatz verwechselt werden darf.
  \item Die angegebenen Macro-F1-/ROC-AUC-Werte beschreiben Klassifikatoren,
        die auf den jeweiligen Datenvarianten trainiert und auf dem
        Original-Testdatensatz geprüft wurden. Das ist kein direkter
        statistischer Nachweis der Verteilungsähnlichkeit der Generatoren.
  \item Erklären, wie Originalbeispiele in die Blöcke injiziert werden und
        wie dadurch die tatsächlichen Quellenanteile pro Klasse entstehen.
  \item Begründen, ob die Datenquellenanteile und Modellvergleiche vor der
        finalen Auswahl festgelegt wurden; andernfalls mögliche
        Auswahlverzerrung diskutieren.
\end{itemize}

\subsection{Grafikbedarf}
\begin{itemize}
  \item \textbf{Vorhanden:} \url{evaluation/plots/model_evaluation_metrics.png}
        für Macro-F1, Test Loss, ROC-AUC und die weiteren im Notebook
        dargestellten Kennzahlen.
  \item \textbf{Vorhanden:} \url{evaluation/plots/feature_importance_comparison.png}
        als Vergleich der Merkmalswichtigkeiten; nur verwenden, wenn die
        Interpretation für die Datengenerierung relevant ist.
  \item \textbf{Ergänzen:} Balken- oder Tabellenansicht der drei
        Datenquellenanteile je Klasse (TabGAN, DAE, Original).
\end{itemize}
